%% ma: L10-B14-0063
%% lop: 10
%% chuong: 5
%% bai: 14
%% dang: 10.14.8
%% loai: TN4
%% muc: VD
%% dap_an: "D"
%% nguon: "chuyen-de-cac-so-dac-trung-cua-mau-so-lieu-khong-ghep-nhom-toan-10-chuong-trinh-moi.pdf (GV. Phan Nhật Linh) (P05), Bài 03 Dạng 2 Phần I Câu 7"
%% kien_thuc: "Lớp 10 Bài 13 (số trung bình), Bài 14 (phương sai, độ lệch chuẩn)"
%% trang_thai: chinh-thuc
%% ghi_chu: "Tự duyệt theo yêu cầu giáo viên 22/10/2026."
\begin{ex}
Điểm thi môn Toán lớp 10A của một trường THPT được trình bày ở bảng tần số sau:
\begin{center}
\begin{tabular}{|c|c|c|c|c|c|c|c|}\hline
Điểm thi & $5$ & $6$ & $7$ & $8$ & $9$ & $10$ & \\\hline
Tần số & $7$ & $5$ & $10$ & $12$ & $4$ & $2$ & $n=40$\\\hline
\end{tabular}
\end{center}
Trong các giá trị dưới đây, giá trị nào gần nhất với phương sai của mẫu số liệu?
\choice{$0{,}94$}{$3{,}94$}{$2{,}94$}{\True $1{,}94$}
\loigiai{$\overline x=\dfrac{287}{40}=7{,}175$; $s^2=\dfrac{1}{40}\sum n_i\left(x_i-7{,}175\right)^2\approx1{,}94$. Chọn D.}
\end{ex}
